\SourceLeaf{400}%
\[
\begin{aligned}
\tfrac{\alpha y'+\beta(y''+y)-\eta}{\alpha x'+\beta(x''+x)-\xi}.\tfrac{y''-y}{x''-x}&=-1\\
\tfrac{\alpha y''+\beta(y+y')-\eta}{\alpha x''+\beta(x+x')-\xi}.\tfrac{y-y'}{x-x'}&=-1
\end{aligned}
\]
Daraus ergibt sich leicht durch Elimination
\[
\begin{aligned}
\xi&=\tfrac{(y-y')(y'-y'')(y''-y)(\alpha-\beta)}{y(x''-x')+y'(x-x'')+y''(x'-x)}\\
&\quad+\alpha.\tfrac{xy(x''-x')+x'y'(x-x'')+x''y''(x'-x)}{y(x''-x')+y'(x-x'')+y''(x'-x)}\\
&\quad+\beta.\tfrac{y(x''x''-x'x')+y'(xx-x''x'')+y''(x'x'-xx)}{y(x''-x')+y'(x-x'')+y''(x'-x)}
\end{aligned}
\]
\par
Der Wert von \(\eta\) folgt aus dem Wert von \(\xi\), wenn man darin alle \(x,x',x''\) mit den entsprechenden \(y,y',y''\) vertauscht, wie man auch von vornherein leicht einsehen kann.

Die Koordinaten des Schnittpunkts der ersten und der letzten Senkrechten ergeben sich, wie man leicht sieht, aus \(\xi\) und \(\eta\), wenn man \(x\) mit \(x'\) und \(y\) mit \(y'\) vertauscht. Dabei ändern \(\xi\) und \(\eta\) ihre Werte jedoch nicht, denn die Koordinaten der Punkte \(A,A',A''\) gehen offenbar in beide Ausdrücke auf gleiche Weise ein. Daher fällt dieser zweite Schnittpunkt mit dem ersten zusammen, und aus demselben Grund fällt der dritte Schnittpunkt ebenfalls von selbst mit den beiden ersten zusammen.

Für den Schwerpunkt gilt übrigens offenbar
\[
n=\tfrac23
\]
also
\[
\alpha=\beta=\tfrac13
\]
und daher
\[
\begin{aligned}\xi&=\tfrac13(x+x'+x'')\\\eta&=\tfrac13(y+y'+y'')\end{aligned}
\]
\par
Für den Mittelpunkt des umschriebenen Kreises gilt
\[
n=1
\]
also
\[
\alpha=0,\quad\beta=\tfrac12
\]
\par
Für den Schnittpunkt der Senkrechten aus \(A\) und den übrigen Eckpunkten selbst gilt
\[
n=0
\]
